\section{Preliminaries}


\subsection{Notations}
\label{sec:notation}
We work over a small finite field \(\Field\) of order
\(p=|\Field|\) and its degree-\(\kappa\) extension
\(\ExtensionField=\mathsf{GF}(p^\kappa)\), where \(p\) is a
prime power, not necessarily the characteristic, and
$
\ExtensionDegree{\ExtensionField}{\Field}=\kappa.
$
Thus \(\Field\subset\ExtensionField\).
For a security parameter \(\lambda\), choose \(\kappa\)
to satisfy the soundness bounds below; in particular, require
\(
    \lvert\ExtensionField\rvert
    =
    p^\kappa
    \geq
    2^\lambda.
\)
We use \(N\) to denote the witness size measured in
\(\Field\)-elements and set \(m=N/\kappa\), assuming that
\(\kappa\mid N\). We restrict ourselves
to parameter regimes satisfying
\(
    p
    >
    \max\left\{
        \kappa,\,
        4\left\lceil\log_2 N\right\rceil
    \right\}.
\)
These are field-cardinality conditions, not restrictions on the
characteristic.

The main exposition specializes to odd prime fields.
Section~\ref{sec:binary-packing} gives the binary-field
instantiation using a small extension as the working field.

We use the following notational conventions throughout this paper.
Scalars in \(\Field\) are denoted by lowercase letters,
\eg, \(a,b,c\).
Matrices are denoted by bold uppercase letters,
\eg, \(\bm A,\bm B,\bm C\),
whereas vectors are denoted by bold lowercase letters and are
interpreted as column vectors by default.
For example, \(\bm l^\top\) and \(\bm r\) denote a row vector
and a column vector, respectively.
Matrices, vectors, and scalars over \(\ExtensionField\) are
distinguished by an overline,
\eg, \(\ext{\bm A}\), \(\ext{\bm l}\), and \(\ext{a}\).

Let
\(\overline{\beta}_1,\ldots,\overline{\beta}_\kappa\)
be a basis of \(\ExtensionField\) over \(\Field\), and let
\(\overline{\beta}_1^\dagger,\ldots,
\overline{\beta}_\kappa^\dagger\)
denote its trace-dual basis. 
The polynomial rings over \(\Field\) and \(\ExtensionField\) are
denoted by \(\Field[X]\) and \(\ExtensionField[X]\), respectively.
Formal indeterminates are denoted by uppercase letters
\(X,Y,Z\), whereas sampled random challenges are denoted by
\(\overline{x},\overline{y},\overline{z}\).
Indices are denoted by \(i,j,\ell\).
% Sets are denoted using blackboard bold font,\eg,
% The set of natural numbers is $\mathbb{N}$,
% the set of powers of two is $2^\mathbb{N}$,
% and the set of real numbers is $\mathbb{R}$.
$[n]$ denotes $\{ 1, \ldots, n \}$.
$|\cdot|$ denotes the cardinality.


We also use the following matrix and vector operations:
\begin{itemize}[noitemsep, topsep = 0pt]
\item inner product $c = \langle \bm a, \bm b \rangle$; addition ($\bm C = \bm A + \bm B$); transposition ($\bm C = \bm A^\top$);
\item multiplication
($\bm C = \bm A\bm B$); Kronecker (tensor) product ($\bm C = \bm A \otimes \bm B$);
\item Hadamard (element-wise) product $\bm{A} \circ \bm{B}$;
\item concatenation of vectors $\bm{c} = (\bm c_1 \| \bm c_2 \|\cdots \| \bm c_\kappa)$; and
\item horizontal packing of column vectors $\bm C = [\bm c_1 \mid \bm c_2 \mid \cdots \mid \bm c_\kappa]$.
\end{itemize}
    
In our complexity analysis, we use \(\bb\), \(\be\), and \(\ee\)
to denote the costs of base-field--base-field,
base-field--extension-field, and extension-field--extension-field
multiplications, respectively, and subsume the cost of additions
into these multiplication costs. Under schoolbook extension-field
arithmetic, we have
\(
    \be=\kappa\bb
\) and
\(
    \ee=\kappa\be=\kappa^2\bb
\).


The prover and the verifier are denoted by $\Prover$ and $\Verifier$, respectively. 
A relation is denoted by $\Relation = \{\text{ public input }:\text{ private input }| \text{ constraint }\}$.

\subsection{SNARK Protocol in Matrix Form}
\label{sec:defgeneral}
Many modern SNARK protocols arithmetize the circuit-satisfiability
(\(\mathsf{Circ\text{-}SAT}\)) problem as a low-degree polynomial
identity and use the sum-check protocol to reduce its verification
to checking multilinear-extension (MLE) evaluations at random
challenge points.
Although equivalent to the matrix formulation,
this polynomial view would make our mathematical development in this paper unnecessarily cumbersome.
We therefore formulate the SNARK relation and the associated protocols directly in matrix form.


\begin{definition}[PLONK constraint system for circuit
satisfiability~\cite{iacr/GabizonWC19,iacr/AmbronaSTW22}]
\label{def:plonk}
An arithmetic circuit with \(N\) gates and \(n_{\mathsf{io}}\)
public inputs and outputs is specified by five selector vectors
\(
    \bm q_L,\bm q_R,\bm q_M,\bm q_O,\bm q_C
    \in\Field^N
\)
and an augmented wiring permutation
\(
    \sigma:
    [3N+n_{\mathsf{io}}]
    \to
    [3N+n_{\mathsf{io}}],
\)
which encodes both the wire-copy constraints and the connections
to the public inputs and outputs.
Given a public-I/O vector
\(\bm{\mathsf{io}}\in\Field^{n_{\mathsf{io}}}\), let
\(\bm w_L,\bm w_R,\bm w_O\in\Field^N\) denote the left, right,
and output wire assignments, respectively, and define the
augmented assignment
\(
    \bm w
    \coloneqq
    (\bm w_L\mathbin{\|}\bm w_R\mathbin{\|}\bm w_O
    \mathbin{\|}\bm{\mathsf{io}})
    \in\Field^{3N+n_{\mathsf{io}}}.
\)
The wire assignment satisfies the circuit if:
\begin{align}
    \bm q_L\circ\bm w_L
    +\bm q_R\circ\bm w_R
    +\bm q_M\circ(\bm w_L\circ\bm w_R)
    +\bm q_O\circ\bm w_O
    +\bm q_C
    &=\bm 0_N,
    \label{eq:plonk1}\\
    \sigma(\bm w)
    &=\bm w.
    \label{eq:plonk2}
\end{align}
Here,
\(\sigma(\bm w)\) denotes the coordinate permutation of
\(\bm w\) induced by \(\sigma\).
\end{definition}


To express MLE evaluations directly in
vector and matrix form, we introduce tensor spans.

\begin{definition}[Tensor span]
\label{def:tensorspan}
Let \(n \in 2^\mathbb{N}\) and let
\(
    \bm x
    =
    (x_1,\ldots,x_{\log_2 n})
    \in\Field^{\log_2 n}.
\)
The \emph{tensor span} of \(\bm x\) is the length-\(n\) vector
\(
    \tensorspan{\bm x}
    \coloneqq
    \bigotimes_{i=1}^{\log_2 n} (1-x_i, x_i)^\top
    \in\Field^n.
\)
The same construction applies to
\(\ext{\bm x}\in\ExtensionField^{\log_2 n}\).
We define the sets of tensor spans over \(\Field\) and
\(\ExtensionField\), respectively, by:
\begin{align*}
    \TensorSpace(n)
    &\coloneqq
    \left\{
        \tensorspan{\bm x}
        \;\middle|\;
        \bm x\in\Field^{\log_2 n}
    \right\}
    \subseteq\Field^n,
    \qquad
    \ext{\TensorSpace}(n)
    \coloneqq
    \left\{
        \tensorspan{\ext{\bm x}}
        \;\middle|\;
        \ext{\bm x}\in\ExtensionField^{\log_2 n}
    \right\}
    \subseteq\ExtensionField^n.
\end{align*}
\end{definition}

\begin{lemma}[Schwartz-Zippel-DeMillo-Lipton Lemma]
\label{lemma:schwartz}
Let
\( \phi(\bm{x}) : \Field^\upsilon \to \Field \) be a nonzero \( \upsilon \)-variate polynomial such that the degree of \( \phi \) in each variable is at most \( d \). Let \( \bm{x} \ChooseRandom \mathbb{X}^{\upsilon} \) for some finite set \( \mathbb{X} \subseteq \Field \). Then:
\[
    \Pr_{\bm{x} \ChooseRandom \mathbb{X}^{\upsilon}}\!\left[ \phi(\bm{x}) = 0 \right]
    \le \frac{\upsilon \cdot d}{|\mathbb{X}|}.
\]
\end{lemma}

\begin{corollary}[Zero check in matrix form]
\label{corollary:schwartz}
Let \(n\) be a positive power of two and let
\(
    \bm a\in\Field^n\setminus\{\bm 0_n\}.
\)
Viewing \(\bm a\) as a vector over \(\ExtensionField\), sample
\(
    \ext{\bm x}
    \ChooseRandom
    \ExtensionField^{\log_2 n}.
\)
Then the probability that the MLE of \(\bm a\)
vanishes at \(\ext{\bm x}\) is bounded by:
\[
    \Pr_{\ext{\bm x}\ChooseRandom
        \ExtensionField^{\log_2 n}}
    \!\left[
        \left\langle
            \bm a,
            \tensorspan{\ext{\bm x}}
        \right\rangle
        =0
    \right]
    \leq
    \frac{\log_2 n}{\lvert\ExtensionField\rvert}.
\]
\end{corollary}

\begin{definition}[Inner-product proof]
\label{def:ip}
Let \(m\) be a positive power of two. Suppose that the prover
\(\Prover\) is bound to two vectors
\(\ext{\bm a},\ext{\bm b}\in\ExtensionField^m\), and that the
verifier \(\Verifier\) holds the corresponding commitments.
Given a claimed value \(\ext{v}\in\ExtensionField\), an
\emph{inner-product proof} is an interactive proof protocol for the
relation:
\[
    \Language_{\mathsf{ip}}
    \coloneqq
    \left\{
        \bigl(\ext{v}\in \ExtensionField,~      \ext{\bm a},\ext{\bm b}\in\ExtensionField^m
        \;\middle|\;
        \begin{array}{l}
            \ext{v}
            =\langle\ext{\bm a},\ext{\bm b}\rangle
        \end{array}
    \right\}.
\]
\end{definition}

Inner-product proofs can be constructed using either the
sum-check protocol~\cite{crypto/Thaler13} or folding
techniques~\cite{asiacrypt/CCYY25}. Both approaches yield the
following reduction.

We represent sum-check round polynomials by their coefficients.
The consistency check is \(g(0)+g(1)=c\); it requires no
division by \(2\) and applies in characteristic two as well.

\begin{fact}[\bf Inner-product reduction]
\label{fact:ip}
There exists an inner-product proof whose reduction phase requires
\(4m+O(\log m)\) multiplications in \(\ExtensionField\), communicates
\(2\log_2 m+2\) $\ExtensionField$-elements, and requires
\(O(\log m)\) $\ExtensionField$-multiplications from the verifier.
Over \(\log_2 m\) rounds, the verifier samples challenges collected
in
\(
    \ext{\bm x}\in\ExtensionField^{\log_2 m},
\)
and the protocol reduces the inner-product relation to two
terminal MLE claims:
\[
    \ext{v}_1
    \EqOrNot
    \left\langle
        \ext{\bm a},
        \tensorspan{\ext{\bm x}}
    \right\rangle
    \quad\wedge\quad
    \ext{v}_2
    \EqOrNot
    \left\langle
        \ext{\bm b},
        \tensorspan{\ext{\bm x}}
    \right\rangle,
\]
where \(\ext{v}_1,\ext{v}_2\in\ExtensionField\) are sent by the
prover. The soundness error is at most
\(
    \frac{2\log_2 m}{\lvert\ExtensionField\rvert}.
\)
\end{fact}

\begin{definition}[Multilinear polynomial commitment scheme in matrix form]
\label{def:pcs}
Let \(n \in 2^\mathbb{N}\). A multilinear PCS over \(\Field\), supporting evaluations over
\(\ExtensionField\), consists of the following algorithms $(\Setup, \Com, \Open)$:
\begin{itemize}[leftmargin=*,labelsep=0.5em]
    \item
    \(
        \pp\leftarrow\Setup(1^\lambda,n)
    \)
    takes the security parameter and the committed vector length
    \(n\), and outputs public parameters \(\pp\).

    \item
    \(
        C\leftarrow\Com(\pp,\bm a;\rho)
    \)
    takes a vector \(\bm a\in\Field^n\) and fresh randomness
    \(\rho\), and outputs a commitment \(C\). We write
    \(\Com(\bm a;\rho)\) when \(\pp\) is clear from context.

    \item The opening protocol
    \(
        \mathsf{bit}
        \leftarrow
        \Open\!\left(
            \pp,C,\ext{v},\ext{\bm x}
            :
            \bm a,\rho
        \right)
    \)
    is an interactive argument between \(\Prover\) and
    \(\Verifier\) for the relation:
    \[
        \Relation_{\mathsf{open}}
        \coloneqq
        \left\{
            \bigl(
                \pp,C,\ext{v},\ext{\bm x}
            \bigr)
                :
            \bigl(
                \bm a,\rho
            \bigr)
            \;\middle|\;
            \begin{aligned}
                C=\Com(\pp,\bm a;\rho) \wedge
                \ext{v}=\left\langle
            \bm a,
            \tensorspan{\ext{\bm x}}
        \right\rangle
            \end{aligned}
        \right\}.
    \]
    At the end of the protocol, \(\Verifier\) outputs
    \(\mathsf{bit}\in\{\TRUE,\FALSE\}\).
\end{itemize}
\end{definition}



\subsection{Extension-Field Arithmetic}

\begin{fact}[Base-Field Embedding]
\label{fact:injective}
Let\/ \(\ExtensionField/\Field\) be an extension.
The canonical inclusion
\(
    \iota:\Field\hookrightarrow\ExtensionField
\)
is injective and induces an injective inclusion
\(
    \Field[X]\hookrightarrow\ExtensionField[X].
\)
\end{fact}

Conversely, the field trace defines a map from \(\ExtensionField\) to \(\Field\).
\begin{definition}[Extension-field trace]
\label{def:trace}
The \emph{field trace} from \(\ExtensionField = \mathsf{GF}(p^\kappa)\) to \(\Field = \mathsf{GF}(p)\)
is the \(\Field\)-linear map:
\[
    \trace:\ExtensionField\to\Field,
    \qquad
    \ext{\alpha}
    \longmapsto
    \trace(\ext{\alpha})
    \coloneqq \ext{\alpha} + \ext{\alpha}^p + \ext{\alpha}^{p^2} + \cdots + \ext{\alpha}^{p^{\kappa-1}}.
\]
In particular, for all \(a,b\in\Field\) and
\(\ext{\alpha},\ext{\beta}\in\ExtensionField\),
\[
    \trace\!\left(
        a\ext{\alpha}+b\ext{\beta}
    \right)
    =
    a\trace(\ext{\alpha})
    +
    b\trace(\ext{\beta}).
\]
\end{definition}

\begin{definition}[Basis and trace-dual basis]
\label{def:basis}
An ordered tuple
\(
    \mathcal B
    =
    \bigl(
        \ext{\beta}_1,\ldots,\ext{\beta}_\kappa
    \bigr)
    \in\ExtensionField^\kappa
\)
is a \emph{basis} of \(\ExtensionField\) over \(\Field\) if every
\(\ext{\alpha}\in\ExtensionField\) admits a unique representation:
\[
    \ext{\alpha}
    =
    \sum_{i=1}^{\kappa}a_i\ext{\beta}_i,
    \qquad
    a_1,\ldots,a_\kappa\in\Field.
\]
The \emph{trace-dual basis} of \(\mathcal B\) is the unique basis
\(
    \mathcal B^\dagger
    =
    \bigl(
        \ext{\beta}_1^\dagger,\ldots,
        \ext{\beta}_\kappa^\dagger
    \bigr)
\)
satisfying:
\[
    \trace\!\left(
        \ext{\beta}_i\ext{\beta}_j^\dagger
    \right)
    =
    \delta_{ij}
    \qquad
    \text{for all }i,j\in[\kappa],
\]
where \(\delta_{ij}\) is the Kronecker delta. Consequently, the
coordinates of \(\ext{\alpha}\) with respect to \(\mathcal B\)
can be recovered as:
\[
    a_i
    =
    \trace\!\left(
        \ext{\alpha}\ext{\beta}_i^\dagger
    \right),
    \qquad
    i\in[\kappa].
\]
\end{definition}

Moreover, when the interpolation domain lies in \(\Field\), the Lagrange basis for degree-\(<\kappa\) polynomials over \(\ExtensionField\) consists entirely of polynomials in \(\Field[X]\).


\begin{definition}[Lagrange basis over the extension field]
\label{def:lagrange}
Assume that \(p\geq\kappa\), and fix an ordered set of distinct
nodes \(\mathcal H=\{\alpha_1,\ldots,\alpha_\kappa\}\subseteq\Field\).
The indices \(j\in[\kappa]\) label these nodes; they are not
identified with integer multiples of \(1\) in the field.
The vanishing polynomial of \(\mathcal H\) is:
\[
    \Psi(X)
    \coloneqq
    \prod_{j=1}^{\kappa}(X-\alpha_j)
    \in\Field[X].
\]
For each \(j\in[\kappa]\), define the Lagrange basis polynomial
\(
    \psi_j(X)
    \coloneqq
    \prod_{\substack{\ell\in[\kappa]\\ \ell\neq j}}
    \frac{X-\alpha_\ell}{\alpha_j-\alpha_\ell}
    \in\Field[X].
\)
These polynomials satisfy
\[
    \psi_j(\alpha_i)=\delta_{ij}
    \qquad
    \text{for all }i,j\in[\kappa].
\]
Moreover, \(\{\psi_j(X)\}_{j=1}^{\kappa}\) forms an
\(\ExtensionField\)-basis for the space of polynomials in
\(\ExtensionField[X]\) of degree less than \(\kappa\).
Consequently, every \(g(X)\in\ExtensionField[X]\) with
\(\deg g<\kappa\) admits the unique representation
\[
    g(X)
    =
    \sum_{j=1}^{\kappa}
    \ext{a}_j\psi_j(X),
    \qquad
    \ext{a}_j=g(\alpha_j)\in\ExtensionField.
\]
\end{definition}
